\documentclass[11pt]{article}
\usepackage[margin=2.5cm]{geometry}
\usepackage{amsmath,amssymb}
\usepackage{bm}
\usepackage{algorithm}
\usepackage{algorithmic}

\newcommand{\avg}[1]{\langle #1 \rangle}
\newcommand{\dd}{\mathrm{d}}
\newcommand{\dt}{\Delta t}

\title{Non-stationary memory kernel extraction\\via multi-origin least squares}
\author{}
\date{}

\begin{document}
\maketitle

%% ============================================================
\section{Generalized Langevin equation}
%% ============================================================

We consider a collective variable $q$ whose dynamics obey the generalized Langevin equation (GLE)
\begin{equation}\label{eq:gle}
  \dot v(t) = F\bigl(q(t)\bigr)
    - \int_0^{t} K(t-s)\,v(s)\,\dd s
    + \eta(t),
\end{equation}
where $v = \dot q$, $F(q)$ is the mean force (potential of mean force derivative),
$K(t)$ is the memory kernel, and $\eta(t)$ is the projected noise.
The integral runs from the initial time $t=0$ of the trajectory.

%% ============================================================
\section{Non-stationary Volterra equation}
%% ============================================================

Multiplying Eq.~\eqref{eq:gle} evaluated at time $t_0 + \tau$ by $v(t_0)$
and taking the ensemble average over independent trajectories yields
\begin{equation}\label{eq:volterra_full}
  \avg{v(t_0)\,\dot v(t_0+\tau)}
  = \avg{v(t_0)\,F\bigl(q(t_0+\tau)\bigr)}
  - \int_0^{t_0+\tau} K(t_0+\tau-s)\,\avg{v(t_0)\,v(s)}\,\dd s
  + \avg{v(t_0)\,\eta(t_0+\tau)}.
\end{equation}
For a cold start of the auxiliary (bath) variables ($s(0)=0$) and recording
the deterministic acceleration $a_{\mathrm{det}} = F - \gamma s$ without the
stochastic noise, the noise--velocity correlation vanishes:
$\avg{v(t_0)\,\eta(t_0+\tau)} = 0$ for all $t_0 \geq 0$.

Defining the orthogonal correlation
\begin{equation}
  C_{\perp}(t_0,\tau)
  \;=\; \avg{v(t_0)\bigl[a_{\mathrm{det}}(t_0+\tau) - F(q(t_0+\tau))\bigr]},
\end{equation}
we obtain the \textbf{exact} relation
\begin{equation}\label{eq:volterra_exact}
  \boxed{
  C_{\perp}(t_0,\tau)
  = -\int_0^{t_0+\tau} K(s)\;\avg{v(t_0)\,v(t_0+\tau-s)}\;\dd s.
  }
\end{equation}
Note that the integral runs from $0$ to $t_0+\tau$, \emph{not} from $0$ to $\tau$.
This is the key point: the memory of the system extends back to the true initial time.

%% ============================================================
\section{Standard approach: single time origin}
%% ============================================================

At $t_0 = 0$, Eq.~\eqref{eq:volterra_exact} reduces to the standard
Volterra equation of the first kind:
\begin{equation}\label{eq:volterra_t0}
  C_{\perp}(0,\tau)
  = -\int_0^{\tau} K(s)\;C_{vv}(\tau-s)\;\dd s,
\end{equation}
where $C_{vv}(\tau) = \avg{v(0)\,v(\tau)}$ is the non-stationary velocity
autocorrelation.  This can be solved sequentially (rectangular or trapezoidal
rule) or by differentiating to obtain a better-conditioned second-kind equation.

\paragraph{Limitation.}
Only data at $t_0 = 0$ is used.  With $N$ independent trajectories,
the statistical noise in the correlation functions scales as $N^{-1/2}$.
Using additional time origins $t_0 > 0$ could reduce the variance, but
naively slicing the trajectory and applying Eq.~\eqref{eq:volterra_t0}
from the new ``$t=0$'' ignores the pre-memory integral from $\tau$ to
$t_0+\tau$, yielding incorrect kernels.

%% ============================================================
\section{Multi-origin least-squares formulation}
%% ============================================================

We discretize time as $t_n = n\,\dt$ and write the kernel as
$K_j = K(j\,\dt)$ for $j = 0, 1, \ldots, n_K - 1$.
Eq.~\eqref{eq:volterra_exact} at origin $t_0$ and lag $\tau$ becomes
\begin{equation}\label{eq:discrete}
  G_{t_0,\tau}
  = -\dt \sum_{j=0}^{j_{\max}-1} K_j \; A_{t_0,\tau,j},
\end{equation}
where
\begin{align}
  G_{t_0,\tau}
    &= \avg{v(t_0)\;\bigl[a_{\mathrm{det}}(t_0{+}\tau) - F(q(t_0{+}\tau))\bigr]},
    \label{eq:G}\\[4pt]
  A_{t_0,\tau,j}
    &= \avg{v(t_0)\;v(t_0{+}\tau{-}j\,\dt)},
    \label{eq:A}\\[4pt]
  j_{\max}
    &= \min\!\bigl(n_K,\; t_0{+}\tau{+}1\bigr).
    \label{eq:jmax}
\end{align}
Crucially, $j$ ranges up to $t_0 + \tau$ (capped at $n_K$), so the velocity
correlations reach back to $v(0)$---the full pre-memory is included.

\subsection{Building the linear system}

For each time origin $t_0 = 0, \dt, \ldots, t_0^{\max}$ and each lag
$\tau = \dt, 2\dt, \ldots, \tau^{\max}$ (subject to $t_0 + \tau < n_K$),
Eq.~\eqref{eq:discrete} provides one linear equation in the unknowns
$\bm{K} = (K_0, K_1, \ldots, K_{n_K-1})^\top$.

Stacking all equations:
\begin{equation}\label{eq:linsys}
  \bm{G} = -\dt\;\bm{A}\;\bm{K},
\end{equation}
where $\bm{G} \in \mathbb{R}^M$ is the vector of measured left-hand sides,
$\bm{A} \in \mathbb{R}^{M \times n_K}$ is the matrix of velocity correlations,
and $M \gg n_K$ (overdetermined system).

\subsection{Smoothness-regularized solution}

Adjacent columns of $\bm{A}$ are highly correlated: column~$j$ contains
$\avg{v(t_0)\,v(t_0{+}\tau{-}j\,\dt)}$ while column~$j{+}1$ contains
the same correlation shifted by one timestep.
Standard (zeroth-order) Tikhonov regularization $\lambda\|\bm{K}\|^2$
penalizes the \emph{magnitude} of the kernel, systematically shrinking
$K(0)$ toward zero.

Instead, we use first-order Tikhonov regularization, penalizing the
roughness of the kernel:
\begin{equation}
  \bm{K}^*
  = \arg\min_{\bm{K}}
    \Bigl\{
      \bigl\|\bm{A}\bm{K} + \bm{G}/\dt\bigr\|^2
      + \lambda\,\|\bm{D}\bm{K}\|^2
    \Bigr\},
\end{equation}
where $\bm{D}$ is the $(n_K{-}1) \times n_K$ first-difference matrix,
$(\bm{D}\bm{K})_i = K_{i+1} - K_i$.
This encourages a smooth kernel without shrinking its magnitude.
The closed-form solution is
\begin{equation}\label{eq:solution}
  \boxed{
  \bm{K}^*
  = \bigl(\bm{A}^\top\bm{A} + \lambda\,\bm{D}^\top\bm{D}\bigr)^{-1}
    \bm{A}^\top\bigl(-\bm{G}/\dt\bigr).
  }
\end{equation}

\paragraph{Automatic parameter selection.}
The regularization parameter $\lambda$ is selected by L-curve corner
detection: the curve $\bigl(\log\|\bm{A}\bm{K}_\lambda + \bm{G}/\dt\|^2,\;
\log\|\bm{D}\bm{K}_\lambda\|^2\bigr)$ is computed on a log-spaced grid
of~$\lambda$ values, and the point of maximum curvature is chosen.
This balances the data fidelity (residual) against the smoothness penalty
without requiring knowledge of the noise level.

%% ============================================================
\section{Why the standard Volterra at $t_0 > 0$ fails}
%% ============================================================

Slicing the trajectory at frame $t_0$ and applying the standard Volterra
solver implicitly assumes
\begin{equation}\label{eq:wrong}
  C_{\perp}(t_0,\tau)
  \stackrel{?}{=} -\int_0^{\tau} K(s)\;\avg{v(t_0)\,v(t_0+\tau-s)}\;\dd s.
\end{equation}
Comparing with the exact relation~\eqref{eq:volterra_exact}, the missing
\emph{pre-memory} term is
\begin{equation}
  \mathcal{P}(t_0,\tau)
  = -\int_{\tau}^{t_0+\tau} K(s)\;\avg{v(t_0)\,v(t_0+\tau-s)}\;\dd s,
\end{equation}
which involves $\avg{v(t_0)\,v(k)}$ for $k < t_0$---the velocity
correlations reaching \emph{before} the time origin.  This term is generally
nonzero and, if neglected, corrupts the extracted kernel.

The LSQ formulation avoids this entirely: by setting
$j_{\max} = \min(n_K, t_0+\tau+1)$, each row of $\bm{A}$ automatically
includes the pre-memory correlations.

%% ============================================================
\section{Algorithm summary}
%% ============================================================

\begin{algorithm}[H]
\caption{Multi-origin least-squares kernel extraction}
\begin{algorithmic}[1]
\REQUIRE Trajectories $\{q^{(n)}(t),\, v^{(n)}(t),\, a_{\mathrm{det}}^{(n)}(t)\}_{n=1}^N$,
         mean force $F(q)$, timestep $\dt$, kernel length $n_K$,
         max origin $t_0^{\max}$, max lag $\tau^{\max}$
\STATE Compute $a_{\perp}^{(n)}(t) = a_{\mathrm{det}}^{(n)}(t) - F\bigl(q^{(n)}(t)\bigr)$
       for all $n$, $t$
\FOR{$t_0 = 0,\, \dt,\, \ldots,\, t_0^{\max}$}
  \STATE $\tilde v_{t_0}^{(n)} = v^{(n)}(t_0) - \overline{v(t_0)}$
         \hfill\COMMENT{center initial velocity}
  \STATE $\tau_{\mathrm{end}} = \min(\tau^{\max},\; n_K - t_0 - 1)$
  \FOR{$\tau = \dt,\, 2\dt,\, \ldots,\, \tau_{\mathrm{end}}$}
    \STATE $G \leftarrow \frac{1}{N}\sum_n \tilde v_{t_0}^{(n)}\; a_{\perp}^{(n)}(t_0+\tau)$
    \STATE $j_{\max} \leftarrow \min(n_K,\; t_0 + \tau + 1)$
    \FOR{$j = 0, 1, \ldots, j_{\max}-1$}
      \STATE $A_j \leftarrow \frac{1}{N}\sum_n \tilde v_{t_0}^{(n)}\; v^{(n)}(t_0+\tau - j\,\dt)$
    \ENDFOR
    \STATE $A_0 \leftarrow A_0/2$;\; $A_{j_{\max}-1} \leftarrow A_{j_{\max}-1}/2$
           \hfill\COMMENT{trapezoidal rule}
    \STATE Append row $(G,\; \dt\,A_0,\; \dt\,A_1,\;\ldots)$ to the linear system
  \ENDFOR
\ENDFOR
\STATE Select $\lambda$ by L-curve corner detection
\STATE Solve $\bm{K}^* = (\bm{A}^\top\bm{A} + \lambda\,\bm{D}^\top\bm{D})^{-1}\bm{A}^\top(-\bm{G})$
\RETURN $K_j^* = K(j\,\dt)$ for $j = 0, \ldots, n_K - 1$
\end{algorithmic}
\end{algorithm}

%% ============================================================
\section{Comparison of approaches}
%% ============================================================

\begin{center}
\begin{tabular}{l c c c}
\hline
Method & \# equations & Pre-memory & Error propagation \\
\hline
Volterra $t_0=0$ & $\sim\tau^{\max}$ & not needed & sequential \\
Avg.\ Volterra $t_0=0\ldots T$ & $\sim\tau^{\max}$ per $t_0$ & needs $K$ estimate & accumulates \\
LSQ $t_0=0\ldots T$ & $\sim T \times \tau^{\max}$ & built-in & none (single solve) \\
\hline
\end{tabular}
\end{center}

%% ============================================================
\section{Validation results}
%% ============================================================

We validate on the 1D harmonic GLE with exponential kernel
$K(t) = \gamma\,e^{-t/\tau}$, $\gamma = 10$, $\tau = 0.1$,
$k = -4$, $k_BT = 1$, $\dt = 0.001$, using $N = 10{,}000$
cold-start trajectories.

\subsection{Equilibrium accuracy}

\begin{center}
\begin{tabular}{l c c}
\hline
Method & Relative $L^2$ error & $K(0)$ (true: 10.0) \\
\hline
Volterra $t_0=0$           & 0.4\% & 10.0 \\
LSQ $t_0=0$                & 0.5\% & 9.9 \\
LSQ $t_0 \leq 0.1$         & 0.6\% & 9.9 \\
\hline
\end{tabular}
\end{center}

All methods recover the kernel accurately at the sub-percent level.
The remaining error is dominated by the time discretization $O(\dt)$.
The cumulative integral $\int_0^t K(s)\,\dd s$ (friction coefficient) is
recovered accurately by all methods.

\subsection{Out-of-equilibrium independence}

The kernel $K(t)$ is a bath property and must be independent of the initial
displacement~$x_0$.  We verify this by extracting $K$ from trajectories
starting at $x_0 = 0$ (equilibrium) and $x_0 = 2$ (displaced by $4\sigma$
from the harmonic minimum).
With the LSQ multi-$t_0$ formulation, the relative difference between the
two extracted kernels is $< 10^{-12}$, confirming exact independence.

\subsection{Requirements}

The method requires \textbf{cold-start auxiliary variables} ($s(0) = 0$
in the OU model for colored noise).
This ensures that $a_{\mathrm{det}} - F = -\gamma\,s$ equals the memory
integral from time zero.
If the auxiliary starts in equilibrium ($s(0) \neq 0$), an unaccounted
transient $\gamma\,s(0)\,e^{-t/\tau}$ corrupts the extraction.

In practice, this corresponds to launching fresh MD trajectories
from a configuration where the bath has been equilibrated at the
restrained CV position.
The memory integral then naturally starts from zero at the trajectory
start, which is the cold-start condition.

%% ============================================================
\section{Regularization and quadrature}
%% ============================================================

Two improvements over a naive implementation are critical for accuracy:

\paragraph{Smoothness penalty.}
Standard (zeroth-order) Tikhonov regularization $\lambda\|\bm{K}\|^2$
penalizes the magnitude of the kernel, systematically shrinking $K(0)$
toward zero.  First-order regularization $\lambda\|\bm{D}\bm{K}\|^2$
penalizes roughness instead, which is physically appropriate for smooth
memory kernels.

\paragraph{Trapezoidal quadrature.}
Adjacent columns of $\bm{A}$ are nearly collinear because
$\avg{v(t_0)\,v(t_0{+}\tau{-}j\dt)}$ and
$\avg{v(t_0)\,v(t_0{+}\tau{-}(j{+}1)\dt)}$ differ by only one timestep
in the lag.  Using the trapezoidal rule (half-weight at the integration
endpoints $j=0$ and $j=j_{\max}{-}1$) breaks this degeneracy and allows
$K_0$ to be properly determined.

\begin{center}
\begin{tabular}{l l c c}
\hline
Quadrature & Penalty & Error & $K(0)$ \\
\hline
Rectangular & $\lambda\|\bm{K}\|^2$       & 7.9\% & 5.5 \\
Rectangular & $\lambda\|\bm{D}\bm{K}\|^2$ & 4.3\% & 8.5 \\
Trapezoidal & $\lambda\|\bm{D}\bm{K}\|^2$ & 0.6\% & 9.9 \\
\hline
\end{tabular}
\end{center}

\end{document}
